\documentclass{article}
\usepackage{import}


\title{Advanced Algorithms\\Homework 03}
\author{Felix Céard-Falkenberg\\7174020}

\usepackage{../homework} % See homework.sty %

\begin{document}

\section{Question 3}
\subsection{\textbf{Give a polynomial-time algorithm to find a schedule of minimum cost by modelling it as a minimum cost flow problem}}

% - m machines
%     - each can process one job at a time
%     - different speed $\ell_i \in \nn$

% - n jobs
%     - same size
%         - processing time is speed $\ell_i$
%     - different levels of urgency
%     - cost function $c_j(t)$
%         - cost of completing job $j$ at time $t$.
%         - non-negative and monotone
    

We model this problem by modeling it as a minimum cost flow problem.
We first describe how we model the problem as a graph problem, and then we prove the polynomial-time complexity of the algorithm.

\textbf{Idea.}
The idea of our algorithm is to model the sequential processing of the jobs by each machine as a multi-source single-sink flow problem, where each machine is a source and has to travers a set of jobs to reach the sink. By adaptively removing edges from the graph and changing the capacity of specific edges, we can find a schedule of minimum cost.

\textbf{Construction of the Graph.}
Let us first create a node for each job $j \in \set{1, \dots, n}$. We will denote each node referencing a job as $J_j$, where $j$ denotes the index of the job. 

We assume that we observe each job from the start, and that no constraints on the jobs exist, such that each machine can process an arbitrary sequence of jobs; Incorporating constraints on the order and availability of the jobs could be easily implemented in this step.

Therefore, we add a connection of each node $J_j$ to each other node $J_i$ with $J_j\neq J_i$ with capacity $\delta\in \nn$. We will comment on the value of $\delta$ later. Naturally, the edge is undirected, or equivalently, two directed edges in opposite directions. This can easily be done in $\mathcal{O}(n^2)$.

Now, we add a node for each of the $m$ machines in the graph, and denote, analogously to the previous paragraph, each machine node as $M_i$. We then connect each machine node $M_i$ to each job node $J_j$ for all $i \in \set{1, \dots, m}$ and $j\in \set{1, \dots, n}$ with capacity $\ell_i$. This can easily be done in $\mathcal{O}(m\cdot n)\leq \mathcal{O}(n^2)$.

Finally, we add a sink node $S$ to the graph. We connect each job node $J_j$ to the sink node $S$ with capacity $\delta$. This can obviously be done in $\mathcal{O}(n)$.

In total, the construction of the graph takes $\mathcal{O}(m + n + n^2 + n \cdot m + n) = \mathcal{O}(m + m \cdot n + n^2)$. Assuming that $n > m$, we can simplify this to $\mathcal{O}(n^2)$, however, there might be cases where $n < m$. We refer to this graph as $G=(V, E)$ with vertices $V$ and edges $E$.

\textbf{Algorithm.}

Before introducing the algorithm, we need to formalize how we define the value of $\delta$.
In short, we set $\delta = 1$ and adjust $\ell_i$ such that they are all integers by multiplying them by some joint factor.
% We introduce this variable to handle real-valued processing time $\ell_i$ of the machines.
% If at least one of the machines has a real-valued processing time, we multiply $\ell_i$ by some factor, such that each machine has an integer processing time. We can then denote as $\delta$ the 

We now describe our algorithm.
We first initialize the flow $f$ to be zero at every edge of the graph. This can obviously be done in $\mathcal{O}(|E|)$.

We can then push flow from each of the sources to the sink


\end{document}
